\section{Experimental Results}
\label{sec:results}

\subsection{Experimental Setup}
\label{sec:data}

All experiments are conducted on a machine equipped with an NVIDIA GeForce RTX 5080 GPU (16\,GB) and an AMD Ryzen 9 9900X 12-core CPU (32\,GB), running Windows 11. The model is implemented in Python 3.12 with PyTorch and CUDA 12.8. Training requires approximately two hours per seed over 400 training steps.

Training instances are generated following the parameter distribution shown in Table~\ref{tab:parameter_distributions}. The policy is trained on instances with $M, N, T \in [5,7]$, drawn at random each episode from this distribution.

Instances are partitioned into three disjoint sets: a training set (generated online from the same distribution at each episode), a validation set (separate instances of the same size range, evaluated every ten epochs under greedy decoding for model selection), and a test set of ten independent benchmark replications per configuration (120 instances total). The twelve test configurations, ranging from $(5,5,5)$ to $(70,100,15)$ and shown in Table~\ref{tab:problem_instance_scales}, span four size tiers (Small, Medium, Large, Battlefield) and are never seen during training. Configurations are restricted to $N \geq M$ (targets at least as numerous as weapons), reflecting resource-constrained defense scenarios.

\begin{table}[htbp]
\centering
\small
\caption{Instance parameter distribution. $\mathcal{Z}(a,b)$ denotes discrete uniform on $\{a, \ldots, b\}$. $D_m$ is the number of stages before a weapon can fire again; $D_m=1$ permits every-stage firing.}
\label{tab:parameter_distributions}
\begin{tabular}{l|l}
\hline
\textbf{Attribute} & \textbf{Distribution} \\ \hline
Target Value $V_n$ & $U(1, 10)$ \\
Reload Time $D_m$ & $\mathcal{Z}(1, 2)$ \\
Ammunition $A_m$ & $\mathcal{Z}(1, \lfloor T/2\rfloor{+}1)$ \\
Start Time $t_{n,\text{start}}$ & $\mathcal{Z}(0, T{-}1)$ \\
End Time $t_{n,\text{end}}$ & $T{-}1$ \\
Damage Rate $P$ & $U(0.3, 0.9)$ \\ \hline
\end{tabular}
\end{table}

\begin{table}[htbp]
\centering
\small
\caption{Test problem configurations. Each configuration uses ten fixed, seed-controlled instances identical across all methods.}
\label{tab:problem_instance_scales}
\begin{tabular}{c|c|c}
\hline
\textbf{Tier} & \textbf{(M,N,T)} & \textbf{Instances} \\ \hline
\multirow{3}{*}{Small}      & (5,5,5)    & 10 \\
                            & (5,7,5)    & 10 \\
                            & (10,15,5)  & 10 \\ \hline
\multirow{3}{*}{Medium}     & (15,15,5)  & 10 \\
                            & (15,20,5)  & 10 \\
                            & (20,30,5)  & 10 \\ \hline
\multirow{3}{*}{Large}      & (30,30,10) & 10 \\
                            & (30,40,10) & 10 \\
                            & (40,50,10) & 10 \\ \hline
\multirow{3}{*}{Battlefield} & (50,50,15) & 10 \\
                            & (50,70,15) & 10 \\
                            & (70,100,15) & 10 \\ \hline
\end{tabular}
\end{table}

%% RESOLVED 2026-08-23: Table~\ref{tab:problem_instance_scales} now states 10 instances/config,
%% matching what every reported run actually used (seed-controlled, fixed across methods).
%% 2026-09-14: distribution table now matches generate_moderate_temporal_instance
%% (V = U(2,6) + U(2,5)*t_start/(T-1), ammo Z(1,floor(T/2)+1), P ~ U(0.3,0.9)).
%% Reload is reported as the firing cadence D_m = prep + 1, so code prep in {0,1}
%% is tabulated as Z(1,2); see DWTA_Simulator.py:317.


Model and training hyperparameters follow standard practice in GNN-based
combinatorial optimization, detailed in Table~\ref{tab:hyperparameters}. No
systematic tuning was performed except on the search budget $K$
(Section~\ref{sec:budget}), whose impact is studied directly. Other configurations may yield better results.

\begin{table}[h]
\centering
\small
\caption{Selected model and training hyperparameters.}
\label{tab:hyperparameters}
\begin{tabular}{ll}
\hline
\textbf{Parameter} & \textbf{Value} \\ \hline
Message-passing layers $L$ & 3 \\
Hidden/embedding dimension $H$ & 128 \\
Attention heads (ablation only, Eq.~\eqref{eq:comm-attn}) & 8 \\
Dropout (residual blocks) & 0.1 \\
Optimizer & Adam \\
Learning rate & $5 \times 10^{-4}$ (cosine decay to $1\%$) \\
Weight decay & $1 \times 10^{-2}$ \\
Gradient-norm clip & 1.0 \\
Instances per training step & 15 \\
Parallel rollouts per instance $P$ & 10 \\
Ammunition bonus $\lambda_{\text{ammo}}$ (Eq.~\eqref{eq:return}) & 0.1 \\
Search budget in training $K_{\text{train}}$ & 3 \\
Search budget at evaluation $K$ & 10 (Table~\ref{tab:budget}) \\ \hline
\end{tabular}
\end{table}

%% Values above verified against common/Dynamic_HYPER_PARAMETER.py and common/DWTA_GNN_comm_sinkhorn.py
%% on 2026-08-23 (EMBEDDING_DIM=128, HEAD_NUM=8, ACTOR_LEARNING_RATE=5e-4, ACTOR_WEIGHT_DECAY=1e-2,
%% ResidualBlock/MultiheadAttention dropout=0.1, Adam). Training-loop values re-verified
%% 2026-09-12 against rl/DWTA_GNN_TRAIN_search_in_loop.py (k_edits=3, k_eval=10, batch 15,
%% para 10, ammo_coef=0.1).
%%
%% RESOLVED 2026-08-23: sinkhorn_scale ablation run (eval_sinkhorn_ablation.py, result/sinkhorn_ablation.txt):
%% scale as trained (0.0213) vs forced to 0, all 12 configs, same instances -- mean 0.1528 vs 0.1527,
%% per-config deltas within +-0.0024 (both signs, noise). Term is INACTIVE; disclosed in one sentence
%% at the end of Section~\ref{sec:comm}.


\subsubsection{Baseline Methods}

We compare against three families of baselines: an exact solver, two
non-learned heuristics, and three learned construction policies.

\paragraph{Exact} The DWTA formulation of Section~\ref{sec:problem} is a
mixed-integer nonlinear program, which we solve with
SCIP~\cite{bestuzheva2023scip} under a 600\,s time limit per instance. When
optimality is not proven within that budget, SCIP reports its best incumbent.

\paragraph{Heuristics} \textsc{Greedy} assigns, at every stage, the
weapon-target pair with the largest immediate marginal destroyed value,
updates that target's remaining value, and repeats until no legal pair
remains. This is the assignment rule of Section~\ref{sec:auction} applied to
every weapon that can fire, which makes it exactly the $\textsc{Myopic}$ rule
of Proposition~\ref{prop:myopic-gap}. It is the baseline that isolates what
the learned fire-or-hold decision contributes: it and our method assign
targets identically and differ only in which weapons fire.
\textsc{Auction} instead matches weapons to targets by the price mechanism of
Bertsekas~\cite{bertsekas1988auction}. Targets carry prices that rise as
weapons bid on them, an incumbent is displaced when outbid, and a weapon
holds fire when no target's remaining value exceeds its price. Prices admit
one weapon per target, so this baseline cannot concentrate several weapons on
a single target.

\paragraph{Learned construction policies} All three build a schedule
autoregressively, one decision at a time, re-observing the state after each.
They differ in how a decision is parameterized and in what the encoder sees.

\textsc{Sequential} decodes one assignment at a time. At
every step it forms a single softmax over all $M \times N$ weapon-target edges
plus one global no-op, samples one edge, applies it, and re-runs the full
edge-aware GNN encoder on the updated state. It therefore chooses \emph{which
weapon acts} and \emph{which target it engages} jointly at each step, and the
kill probability $P_{mn}$ enters the encoder directly as an edge feature.

AM~\cite{kool2018attention} is the attention model of Kool et al., run
unmodified through RL4CO~\cite{berto2023rl4co}. Its encoder is a Transformer
over target nodes, evaluated once per instance; the decoder then attends to
those fixed embeddings through a context query. Because target values change
as damage accumulates, the static-embedding default is unsound here, so
remaining value is supplied to the decoder as an explicit dynamic embedding.
The decision order is fixed in advance, round-major and weapon-minor, so the
step index prescribes which weapon acts and the pointer selects only its
target or a no-op.

POMO~\cite{kwon2020pomo} shares AM's architecture and differs in its rollout
scheme: several trajectories are drawn per instance from forced-distinct
starting actions, and their mean serves as the REINFORCE baseline. We report
POMO under single-path greedy decoding rather than best-of-multistart, so that
its inference budget matches every other method in the table.

All three learned baselines are trained on the same curriculum under the same
optimizer settings, and all neural methods are evaluated zero-shot at every
tier with no fine-tuning on the test configurations.

\subsection{Main Results}
\label{sec:main-results}

Methods are compared on the normalized remaining value, the fraction of the
initial target value still standing when the engagement ends, for which lower
is better.
%% Generated by code_DWTA/eval_final_table.py from the three fire/hold-only
%% checkpoints (seeds 5, 6, 7), result/final_table_nocomm.txt. (NOT
%% final_table_multiseed.txt, which is the older with-communication run and
%% has mean 0.1197 at K=10 rather than the reported 0.1179.)
%% Baseline columns are carried over from the previous benchmark run: verified
%% comparable because both use the same generator, seed 123, and the same ten
%% instances per configuration -- the Sequential mean reproduces to 0.1432 in
%% both, to four decimals.

\begin{table*}[t]
\centering
\scriptsize
\caption{Main results across twelve held-out test configurations.
Normalized remaining value (lower is better). Our column reports the mean over three seeds followed by the across-seed standard deviation. Every method constructs its schedule in a single pass, so the comparison is at matched inference effort; what the local search adds on top is reported separately in Section~\ref{sec:budget}. SCIP uses a 600\,s limit and reports its incumbent; at the two largest configurations it fails to find a competitive solution. Bold marks the best method per configuration.}
\label{tab:main_results}
\begin{tabular}{@{}llcccccccc@{}}
\toprule
Tier & Config & SCIP & Greedy & Auction & AM & POMO & Sequential & Ours \\
\midrule
\multirow{3}{*}{Small}
 & (5,5,5)     & \textbf{0.148} & 0.598 & 0.362 & 0.249 & 0.243 & 0.169 & 0.190 $\pm$ .000 \\
 & (5,7,5)     & \textbf{0.199} & 0.623 & 0.455 & 0.446 & 0.447 & 0.307 & 0.245 $\pm$ .000 \\
 & (10,15,5)   & \textbf{0.192} & 0.619 & 0.418 & 0.446 & 0.488 & 0.300 & 0.225 $\pm$ .001 \\
\midrule
\multirow{3}{*}{Medium}
 & (15,15,5)   & \textbf{0.087} & 0.533 & 0.266 & 0.288 & 0.315 & 0.096 & 0.094 $\pm$ .000 \\
 & (15,20,5)   & \textbf{0.138} & 0.616 & 0.363 & 0.480 & 0.508 & 0.207 & 0.166 $\pm$ .000 \\
 & (20,30,5)   & \textbf{0.176} & 0.638 & 0.448 & 0.548 & 0.584 & 0.295 & 0.209 $\pm$ .000 \\
\midrule
\multirow{3}{*}{Large}
 & (30,30,10)  & 0.015 & 0.355 & 0.078 & 0.164 & 0.093 & 0.016 & \textbf{0.012} $\pm$ .001 \\
 & (30,40,10)  & 0.045 & 0.407 & 0.177 & 0.256 & 0.136 & \textbf{0.032} & 0.037 $\pm$ .002 \\
 & (40,50,10)  & 0.051 & 0.331 & 0.150 & 0.261 & 0.144 & \textbf{0.028} & 0.028 $\pm$ .001 \\
\midrule
\multirow{3}{*}{Battlefield}
 & (50,50,15)  & 0.157 & 0.361 & 0.156 & 0.195 & \textbf{0.077} & 0.088 & 0.084 $\pm$ .002 \\
 & (50,70,15)  & 0.992 & 0.390 & 0.171 & 0.268 & 0.099 & 0.080 & \textbf{0.077} $\pm$ .002 \\
 & (70,100,15) & 0.993 & 0.403 & 0.177 & 0.306 & 0.106 & 0.100 & \textbf{0.096} $\pm$ .001 \\
\midrule
 & \textbf{Mean} & 0.266 & 0.490 & 0.268 & 0.326 & 0.270 & 0.143 & \textbf{0.122} \\
\bottomrule
\end{tabular}
\end{table*}

At matched inference effort, the mean objective across the twelve
configurations is 0.122, against 0.143 for the strongest neural baseline
(Sequential), a 14.9\% relative reduction. Every non-exact baseline is beaten
on the mean.

The table answers two questions at once, and they are the two the
decomposition rests on. \textsc{AM}, \textsc{POMO} and \textsc{Sequential}
all learn the target choice as well as the firing decision, and all three are
beaten: 0.326, 0.270 and 0.143 against our 0.122. Removing target selection
from the policy therefore costs nothing that a learned target head was
supplying. \textsc{Greedy}, at the other end, assigns targets by the same rule
we do and differs only in that no policy decides which weapons fire; it scores
0.490. The firing decision is thus where essentially all of the value sits,
and it is the one part for which Proposition~\ref{prop:myopic-gap} rules out
any per-stage substitute. Learning less does not cost accuracy here; learning
the right part is what buys it.

The per-configuration picture is more informative than the mean, and splits
cleanly by tier. Against SCIP, we lose all six Small and Medium
configurations and win all six Large and Battlefield ones. That is the
expected pattern: at small scale the solver still proves or nearly proves
optimality within its budget, and no heuristic should beat it there; at large
scale its branch-and-bound no longer reaches a competitive incumbent, failing
outright at the two largest configurations. Against Sequential, the
strongest learned baseline, we win 9 of 12 and tie at (40,50,10),
losing only (5,5,5) by 0.021 and (30,40,10) by 0.005. Against all learned
baselines together we are best at 8 of 12; the remaining loss is (50,50,15),
where POMO reaches 0.077 against our 0.084.

All results are reported over three seeds. Across-seed standard deviations
are between 0.000 and 0.002, smaller than the gap to the nearest competing
method at every configuration, so the ordering of methods does not depend on
which training run is used.

%% Measured by code_DWTA/measure_method_timing.py (Greedy, Auction, Sequential,
%% ours) and code_DWTA/measure_rl4co_timing.py (AM, POMO), both on the same ten
%% seed-123 instances per configuration. CUDA-synchronized, warm-up passes
%% discarded, one instance at a time: RL4CO would batch all ten into a single
%% forward, which would credit AM and POMO with parallelism no other method is
%% given. Raw output in result/timing_methods.txt and result/timing_rl4co.txt.

\begin{table}[t]
\centering
\small
\caption{Mean wall-clock time per instance in seconds, averaged over the three
configurations in each tier. Every column is the cost of producing one
schedule: encoding, every network call, the assignment step, and all
environment simulation, measured end to end with CUDA synchronization after
untimed warm-up passes. Our column is the construction pass alone, at
$K{=}0$, which is the budget the table's objectives are reported at; what the
local search adds on top is in Section~\ref{sec:budget}. All methods decode
one instance at a time, since RL4CO would otherwise batch AM and POMO into a
single forward and credit them with parallelism no other method is given.
SCIP is omitted: it is given a fixed 600\,s budget per instance rather than
run to completion, so its time is a setting rather than a measurement
(Appendix~\ref{app:scip}).}
\label{tab:timing}
\begin{tabular}{@{}lcccccc@{}}
\toprule
Tier & Greedy & Auction & AM & POMO & Sequential & Ours \\
\midrule
Small       & \textbf{0.029} & 0.113 & 0.075 & 0.093 & 0.055 & 0.032 \\
Medium      & \textbf{0.053} & 0.159 & 0.153 & 0.153 & 0.122 & 0.063 \\
Large       & \textbf{0.184} & 0.558 & 0.651 & 0.603 & 0.497 & 0.203 \\
Battlefield & \textbf{0.527} & 1.562 & 1.618 & 1.604 & 1.315 & 0.577 \\
\midrule
\textbf{Mean} & \textbf{0.199} & 0.598 & 0.625 & 0.613 & 0.497 & 0.218 \\
\bottomrule
\end{tabular}
\end{table}

Table~\ref{tab:timing} reports what each method costs to run. Our pipeline is
the fastest of every method that learns anything, at 0.218\,s per instance
against 0.497\,s for the sequential decoder and roughly 0.62\,s for AM and
POMO, and the margin widens with scale: at the Battlefield tier it is 0.577\,s
against 1.315\,s. The reason is structural. Construction calls the network
once per stage and decides all $M$ weapons in that one call, whereas a
sequential decoder calls it once per weapon per stage, so its cost grows with
the size of the force while ours does not. Only \textsc{Greedy} is faster, at
0.199\,s, and it has no network to run at all.

Taken together with the objective, this is the result the paper rests on. The
one method that beats us on time is the worst in the table on quality, at
0.490 against our 0.122. The methods that come closest on quality are two to
six times slower. Against \textsc{Sequential}, the strongest learned baseline,
we are better on the objective at nine of twelve configurations and
2.3 times faster on average, so the comparison does not trade one against the
other.

\subsection{Inference-Time Improvement}
\label{sec:budget}

Once a schedule has been constructed it can still be improved without
touching the model. The search flips one fire-or-hold decision, re-simulates
the episode under the same assignment rule, and keeps the change only if the
objective improves, repeating this up to $K$ times
(Section~\ref{sec:local-search}). Nothing is learned here: the budget $K$ buys
computation at inference, not parameters.

Because a flip is kept only on strict improvement, the objective is
non-increasing in $K$, so the informative question is not whether the curve
descends but where it stops. Table~\ref{tab:budget} reports the sweep and, for
each configuration, the mean index of the last accepted flip within a $K{=}30$
probe.

\begin{table}[t]
\centering
\small
\caption{Edit-budget sweep, mean over three seeds. \emph{last accept} is the
mean index of the final accepted flip within a $K{=}30$ probe; a value near 30
indicates the search was still finding improvements when the budget ran out.}
\label{tab:budget}
\begin{tabular}{@{}lccccc@{}}
\toprule
Config & $K{=}0$ & $K{=}3$ & $K{=}10$ & $K{=}30$ & last accept \\
\midrule
(5,5,5)     & 0.1901 & 0.1825 & 0.1740 & 0.1660 & 11.5 \\
(5,7,5)     & 0.2451 & 0.2381 & 0.2347 & 0.2341 & 7.0 \\
(10,15,5)   & 0.2248 & 0.2240 & 0.2212 & 0.2167 & 13.3 \\
(15,15,5)   & 0.0940 & 0.0925 & 0.0918 & 0.0905 & 12.9 \\
(15,20,5)   & 0.1661 & 0.1647 & 0.1597 & 0.1550 & 19.8 \\
(20,30,5)   & 0.2086 & 0.2070 & 0.2040 & 0.1985 & 17.7 \\
(30,30,10)  & 0.0121 & 0.0118 & 0.0113 & 0.0105 & 21.3 \\
(30,40,10)  & 0.0373 & 0.0366 & 0.0355 & 0.0336 & 21.8 \\
(40,50,10)  & 0.0279 & 0.0275 & 0.0263 & 0.0248 & 19.4 \\
(50,50,15)  & 0.0839 & 0.0838 & 0.0837 & 0.0837 & 21.2 \\
(50,70,15)  & 0.0771 & 0.0770 & 0.0769 & 0.0767 & 20.6 \\
(70,100,15) & 0.0960 & 0.0959 & 0.0958 & 0.0956 & 21.7 \\
\midrule
\textbf{Mean} & 0.1219 & 0.1201 & 0.1179 & \textbf{0.1155} & 17.4 \\
\bottomrule
\end{tabular}
\end{table}

Two patterns matter. The gain is modest in absolute terms, 0.1219 to 0.1155,
a 5.2\% relative reduction across the full budget, because the construction
policy is now trained against the searched objective and therefore already
produces schedules the search has little left to fix. That is the intended
behaviour of training through the search, not a weakness of the search.

Second, where the curve flattens scales with problem size. At (5,7,5) the last
accepted flip is at index 7 on average, and the $K{=}10$ and $K{=}30$ columns
differ by 0.001. At the Large and Battlefield configurations the last accepted
flip sits around 20, so the search is still finding improvements when a budget
of 30 expires. A larger schedule offers more slots worth editing, and
because quality never falls as $K$ grows, an operator can
spend accordingly: a tight decision cycle stops at $K{=}3$ and banks most of
the gain, a loose one keeps going.

%% ============================================================
%% COMMENTED OUT 2026-09-14.
%% This section compared two sources of the fire/hold decision: our parallel
%% policy, and the Sequential decoder with its target choices discarded. That
%% second arm is not a clean control. Sequential decides firing and targeting
%% in one step, choosing an edge because a particular target is attractive, so
%% stripping the target away removes the reason the weapon was selected at all.
%% The honest version of this comparison would decide fire/hold weapon by
%% weapon and then assign, against deciding all at once and then assigning,
%% with the same policy in both arms; that needs a separately trained
%% sequential participation policy, since ours is trained to see all weapons
%% at once.
%%
%% It is also redundant. Proposition~\ref{prop:auction-half} holds for ANY
%% firing set, so the guarantee does not care whether the set was produced
%% simultaneously, and Table~\ref{tab:main_results} already shows the parallel
%% pipeline beating the Sequential decoder outright. Restore only if a
%% reviewer asks for the direct comparison AND the sequential participation
%% policy is trained.
%% ============================================================
%% \subsection{Which Decision Needs Learning}
%% \label{sec:firehold-source}
%%
%% Proposition~\ref{prop:auction-half} holds for any firing set, so the source of
%% that set should matter less than its quality. We test this by fixing the
%% pipeline, keeping the same assignment rule, revision, instances and budget, and
%% changing only where the fire/hold decision comes from. The sequential arm runs
%% the sequential decoder to completion on a copy of the live state and keeps only
%% which weapons fired, discarding its target choices exactly as the parallel
%% policy's are discarded.
%%
%% \begin{table}[t]
%% \centering
%% \small
%% \caption{Source of the fire/hold decision, everything else fixed. Normalized
%% remaining value, lower is better. \emph{Seq.\ alone} uses the sequential
%% decoder's own target choices, with no assignment step and no revision.}
%% \label{tab:firehold-source}
%% \begin{tabular}{@{}lcccccc@{}}
%% \toprule
%% & & \multicolumn{2}{c}{parallel} & \multicolumn{2}{c}{sequential} \\
%% \cmidrule(lr){3-4}\cmidrule(lr){5-6}
%% Config & Seq.\ alone & $K{=}0$ & $K{=}10$ & $K{=}0$ & $K{=}10$ \\
%% \midrule
%% (5,5,5)     & 0.1694 & 0.2486 & 0.1794 & 0.1736 & 0.1605 \\
%% (5,7,5)     & 0.3071 & 0.3213 & 0.2696 & 0.3093 & 0.2868 \\
%% (10,15,5)   & 0.2996 & 0.2911 & 0.2640 & 0.3011 & 0.2891 \\
%% (15,15,5)   & 0.0958 & 0.1215 & 0.1108 & 0.1003 & 0.0984 \\
%% (15,20,5)   & 0.2073 & 0.2150 & 0.1831 & 0.2154 & 0.1950 \\
%% (20,30,5)   & 0.2953 & 0.2599 & 0.2302 & 0.3076 & 0.2857 \\
%% (30,30,10)  & 0.0162 & 0.0153 & 0.0122 & 0.0153 & 0.0134 \\
%% (30,40,10)  & 0.0317 & 0.0361 & 0.0331 & 0.0344 & 0.0304 \\
%% (40,50,10)  & 0.0283 & 0.0298 & 0.0287 & 0.0299 & 0.0267 \\
%% (50,50,15)  & 0.0878 & 0.0972 & 0.0917 & 0.0852 & 0.0849 \\
%% (50,70,15)  & 0.0798 & 0.0869 & 0.0811 & 0.0783 & 0.0778 \\
%% (70,100,15) & 0.1003 & 0.1103 & 0.1063 & 0.0992 & 0.0988 \\
%% \midrule
%% \textbf{Mean} & 0.1432 & 0.1528 & \textbf{0.1325} & 0.1458 & 0.1373 \\
%% \bottomrule
%% \end{tabular}
%% \end{table}
%%
%% Unrevised, the sequential source is better (0.1458 vs 0.1528), as
%% Propositions~\ref{prop:par-linear} and~\ref{prop:seq-half} predict. The
%% ordering reverses at $K{=}10$: 0.1325 against 0.1373. Both beat the sequential
%% decoder using its own targets (0.1432), so delegating assignment costs nothing.
%%
%% The reversal is not evidence that simultaneous decisions are better; the
%% $K{=}0$ columns say otherwise. It is that parallel-produced schedules are more
%% improvable at equal budget. The sequential source still wins 7 of 12
%% configurations, so the mean is carried by margin, not by uniform superiority, most visibly at (20,30,5), 0.2302 against 0.2857.
%%
%% One caveat weakens the sequential arm: the revision policy used here was
%% trained on parallel-produced schedules, so the sequential arm is
%% off-distribution. The parallel win is conservative; the sequential numbers are
%% a lower bound.
%%

\subsection{Ablation Studies}
\label{sec:ablation}

Each of the following measures something the framework could have contained
and does not. Multi-agent construction policies normally give agents a way to
see one another before they commit, through a communication layer or a
conflict handler, because agents deciding at the same instant will otherwise
contend for the same resource. Ours has neither. The contended resource here
is the target, no weapon in this pipeline chooses one, and the assignment step
allocates targets after participation is already fixed, so there is nothing
left for the weapons to coordinate over. The same question arises for the
revision step, which could plausibly learn where to edit rather than propose
uniformly at random. In each case we built the richer variant, trained it
under an identical recipe, and kept the smaller one only where the measurement
said the addition did not earn its place.


\subsubsection{Weapon-to-Weapon Communication}
\label{sec:comm-ablation}

Remark~\ref{rem:auction-kills-pileon} predicts that inter-agent communication
should buy nothing here, because the contention it exists to resolve is over
targets and no weapon in this pipeline chooses one. We tested that prediction
by adding one multi-head self-attention layer over the $M$ weapon embeddings,
placed between the encoder and the scoring heads,
\begin{equation}
\label{eq:comm-attn}
\bm{C} = \text{MHA}\!\left(\bm{H}^{\text{wpn}}, \bm{H}^{\text{wpn}}, \bm{H}^{\text{wpn}}\right), \qquad
\bm{h}_m^{\text{wpn,comm}} = \text{Res}\!\left(\text{LN}\!\left(\bm{h}_m^{\text{wpn},(L)} + \bm{C}_m\right)\right),
\end{equation}
with $\bm{H}^{\text{wpn}} = [\bm{h}_1^{\text{wpn},(L)}; \ldots; \bm{h}_M^{\text{wpn},(L)}]$
and $\text{MHA}$ standard self-attention~\cite{vaswani2017attention} over the
weapon axis, and retraining from scratch under an identical recipe: same
three seeds, same curriculum, same budgets, same evaluation protocol.

\begin{table}[t]
\centering
\small
\caption{With and without the weapon-to-weapon self-attention layer. Mean over
three seeds with across-seed standard deviation, twelve held-out
configurations.}
\label{tab:comm-ablation}
\begin{tabular}{@{}lcccc@{}}
\toprule
& $K{=}0$ & $K{=}3$ & $K{=}10$ & $K{=}30$ \\
\midrule
with communication & 0.1242 & 0.1222 & 0.1197 & 0.1169 \\
\textbf{without (reported)} & \textbf{0.1219} & \textbf{0.1201} & \textbf{0.1179} & \textbf{0.1155} \\
\midrule
max.\ across-seed std.\ ($K{=}10$) & \multicolumn{2}{l}{with: 0.0076} & \multicolumn{2}{l}{without: 0.0023} \\
\bottomrule
\end{tabular}
\end{table}

Removing the layer is better at every budget, and the effect is consistent
rather than an artifact of the mean: the no-communication variant is ahead at
8 of 12 configurations and behind at none by more than 0.0015, which is inside
the seed spread. The largest single gap is at (40,50,10), 0.0263 against
0.0359. Across-seed variance also falls by a factor of three.

We take this as confirmation of Remark~\ref{rem:auction-kills-pileon} rather
than a surprise, and note what it does \emph{not} say. It is not evidence that
communication mechanisms are unhelpful in multi-agent construction generally;
it is evidence that in a pipeline where the contended resource is allocated by
a separate mechanism, the coordination those layers provide is redundant. The
prediction and the measurement agree, and the architecture we report is the
smaller one.

\subsubsection{Learned Selection in the Revision Step}
\label{sec:learned-vs-random}

An accept-if-improving rule makes any proposal distribution monotonically
non-degrading (Section~\ref{sec:local-search}). Comparing a learned
proposer against the unrevised base therefore establishes nothing. The
comparison that does is against a random proposer at matched budget.

\begin{table}[t]
\centering
\small
\caption{Learned versus uniformly random slot selection at matched
re-simulation budget, accumulate-and-keep-best operator. Random is averaged
over five repeats.}
\label{tab:learned-vs-random}
\begin{tabular}{@{}lccccc@{}}
\toprule
Config & $K{=}0$ & learned $K{=}3$ & random $K{=}3$ & learned $K{=}10$ & random $K{=}10$ \\
\midrule
(5,5,5)    & 0.2486 & 0.1929 & 0.2085 & 0.1794 & 0.1991 \\
(5,7,5)    & 0.3213 & 0.2764 & 0.2874 & 0.2696 & 0.2773 \\
(10,15,5)  & 0.2911 & 0.2686 & 0.2799 & 0.2640 & 0.2715 \\
(15,15,5)  & 0.1215 & 0.1150 & 0.1169 & 0.1108 & 0.1146 \\
\midrule
\textbf{Mean} & 0.2456 & 0.2132 & 0.2232 & \textbf{0.2060} & 0.2156 \\
\bottomrule
\end{tabular}
\end{table}

The learned selector wins at every configuration, so its contribution is real.
It is also small: random recovers 0.0300 of the 0.0397 total gain. \textbf{76\%
of the improvement comes from re-simulation and the accept rule, not from
learning.}

\subsubsection{The Revision Operator}
\label{sec:operator}

The operator used previously applies every proposed flip and protects only the
reported value. That is a greedy walk, not a neighbourhood search, and
it does not reach a 1-flip local optimum even when run to exhaustion. Accept-if-improving
makes it one.

\begin{table}[t]
\centering
\small
\caption{Accumulate-and-keep-best versus hill climbing at $K{=}10$, matched
budget.}
\label{tab:operator}
\begin{tabular}{@{}lcccc@{}}
\toprule
Config & $K{=}0$ & Accumulate & Hill climb & Hill climb \\
       &         & (learned)  & (learned)  & (random)   \\
\midrule
(5,5,5)   & 0.2486 & 0.1794 & 0.1764 & 0.1622 \\
(5,7,5)   & 0.3213 & 0.2696 & 0.2560 & 0.2323 \\
(10,15,5) & 0.2911 & 0.2640 & 0.2396 & 0.2496 \\
(15,15,5) & 0.1215 & 0.1108 & 0.1049 & 0.1038 \\
\midrule
\textbf{Mean} & 0.2456 & 0.2060 & 0.1942 & \textbf{0.1870} \\
\bottomrule
\end{tabular}
\end{table}

Hill climbing wins at all four configurations and carries the guarantee
accumulation lacks. It also inverts the previous result: under the corrected
operator the learned ordering is \emph{worse} than random (0.1942 vs 0.1870).
The accept test does the work the ordering was doing, and the selector is
off-distribution besides. At this operator arity the learned selector does not
earn its parameters, which is why the final framework has no second network.

%% ============================================================
%% COMMENTED OUT 2026-09-14. Superseded by the timing table in
%% Section~\ref{sec:main-results}.
%%
%% Every quantitative claim below was measured before the assignment step was
%% vectorized. It ran one weapon at a time inside a Python loop, costing
%% O(M^2) small tensor calls per stage; forming the whole marginal matrix at
%% once and updating only the column that changed gives identical actions
%% (verified on 600 random cases) and is 31x faster at 70x100. That reversed
%% the section's conclusion: the pipeline is no longer slower than a
%% sequential pass, it is 2.3x faster on the mean, so the matched-wall-clock
%% argument this section was built to make is no longer needed.
%%
%% Worth salvaging if the cost breakdown is ever wanted again: the split
%% between construction and search, and the point that re-simulation restarts
%% from stage 0 even though the stages before the edited slot are unchanged,
%% which is still true and still an easy speedup.
%% ============================================================
%% \subsection{Inference Cost}
%% \label{sec:efficiency}
%% %% Measured by code_DWTA/measure_final_timing.py: CUDA-synchronized, two
%% %% untimed warm-up instances per configuration drawn from a separate RNG
%% %% stream, ten timed instances, single GPU.
%%
%% \begin{table}[t]
%% \centering
%% \small
%% \caption{Mean per-instance wall-clock time in seconds, tier averages over three
%% configurations each. \emph{Construction} is encode, fire/hold and assignment once
%% per stage; \emph{search} is ten re-simulations; \emph{sequential} is the
%% per-weapon decoding loop.}
%% \label{tab:timing}
%% \begin{tabular}{@{}lcccc c@{}}
%% \toprule
%% Tier & Construction & Search ($K{=}10$) & Total & Sequential & Seq./Constr. \\
%% \midrule
%% Small       & 0.052 & 0.315 & 0.367 & 0.124 & 2.4$\times$ \\
%% Medium      & 0.104 & 0.858 & 0.962 & 0.356 & 3.4$\times$ \\
%% Large       & 0.409 & 3.807 & 4.216 & 1.638 & 4.0$\times$ \\
%% Battlefield & 1.213 & 12.492 & 13.704 & 4.972 & 4.1$\times$ \\
%% \bottomrule
%% \end{tabular}
%% \end{table}
%%
%% Two readings of Table~\ref{tab:timing} matter, and the second corrects a claim
%% it would be easy to make from the first.
%%
%% Construction alone is 2.4 to 4.1 times cheaper than sequential decoding, with
%% the ratio widening as the force grows: construction makes one network call
%% per stage, sequential decoding one per weapon per stage.
%%
%% But construction is not where the time goes. A single search re-simulation
%% costs about as much as an entire construction pass (1.25\,s against 1.21\,s at
%% the largest configuration), even though it invokes no network at all. The
%% assignment step and the simulator, both loop-bound rather than vectorized, dominate
%% everything; the encoder is close to free by comparison. Consequently a budget
%% of $K{=}10$ makes the full pipeline 2.8 times \emph{slower} than a sequential
%% pass, not faster. We state this plainly because the opposite is easy to assume
%% from the decoding-complexity argument and is wrong.
%%
%% The useful comparison is therefore at matched wall clock rather than at
%% matched edit budget. At the largest configuration a sequential pass costs
%% 4.97\,s; construction plus three edits costs 4.96\,s. At that equal budget the
%% objectives are 0.1432 and 0.1222, a 14.7\% relative reduction. This is what
%% the cheap construction pass buys: not a faster answer, but room to search
%% inside the time the alternative needs for a single pass.
%%
%% Both costs are implementation-bound rather than intrinsic. The assignment
%% step is a per-batch Python loop and re-simulation restarts from stage~0 even though
%% stages before the edited slot are unchanged; vectorizing the former and
%% resuming the latter would each cut the dominant term substantially. Neither is
%% done here, so the times above should be read as an upper bound on the cost of
%% this design, not a property of it.
%%

\subsection{Limitations}
\label{sec:limits}

Results are reported over three seeds. Training is non-monotone in steps for every method
we trained here, so checkpoints are selected by best held-out performance
during training rather than final step. We report the full trace, but this
remains a selection procedure. Selection used a cheaper criterion than the
reported table; reported numbers are regenerated for every checkpoint on all
twelve configurations under one protocol.

The non-monotonicity is large enough to state concretely rather than in the
abstract. On one seed the best checkpoint reached 0.1219 at step 260 while the
final step returned 0.1493, and an intermediate evaluation at step 340
degraded to 0.3540 before recovering. On another the best checkpoint was the
final step itself, which suggests the runs are not uniformly converged and
that a longer schedule might still improve them. Reporting the peak alone
would hide all of this. Appendix~\ref{app:traces} therefore gives the complete
evaluation trace for every seed, so that a reader can see what the selection
procedure was choosing from.

One fact tempers the concern. The three selected checkpoints agree closely on
the held-out benchmark despite these divergent trajectories
(Table~\ref{tab:main_results}), so the procedure is selecting reproducibly even
though the path to it is erratic.

We take the instability to be a property of the learning problem rather than
of the implementation. Participation is a binary decision whose reward is
structurally delayed: firing pays immediately, holding pays nothing within the
stage and only later, if at all. Gradient information is correspondingly weak,
since the decision is piecewise constant in the logit. Training through a
fixed search adds a second effect in the same direction: because the search
repairs some of the policy's mistakes before the objective is measured, the
learning signal about those mistakes is attenuated. The search buys a better
objective and a more stable optimization at the cost of a less discriminating
gradient. We report this as a hypothesis consistent with the traces, not as a
measured mechanism.
